\input{preamble.tex}

\title{Отчет по лабораторной работе № 5. \\ Конфигурирование VLAN}
\author{Сущенко Алина Николаевна \\ НПИбд-01-23}
\date{2026}

\begin{document}

\maketitle
\newpage

\tableofcontents

\newpage
\section{Цель работы}
Приобретение практических навыков по настройке VLAN на коммутаторах сети с использованием протоколов VTP и Trunk, а также проверка изоляции трафика между различными VLAN.

\newpage
\section{Выполнение работы}

\subsection{Настройка Trunk-портов на коммутаторах}
В соответствии с заданием и топологией сети необходимо настроить магистральные порты (trunk) на всех коммутаторах для передачи трафика нескольких VLAN.

\begin{enumerate}
\item На центральном коммутаторе \texttt{msk-ansusenko-sw-01} настроили trunk на интерфейсах, соединяющих его с остальными коммутаторами (Рис. \ref{01}):
  \begin{minted}{cisco}
    enable
    configure terminal
    interface g0/1
    switchport mode trunk
    interface g0/2
    switchport mode trunk
  \end{minted}

\begin{center}
    \centering
    \includegraphics[width=\textwidth]{../images/01.png}
    \captionof{figure}{Настройка trunk-портов на msk-ansusenko-sw-01.}
    \label{01}
\end{center}

\item На коммутаторе \texttt{msk-ansusenko-sw-2} настроили trunk на порту \texttt{g0/2} (Рис. \ref{02}):
  \begin{minted}{cisco}
    enable
    configure terminal
    interface g0/2
    switchport mode trunk
    end
    copy running-config startup-config
  \end{minted}

\begin{center}
    \centering
    \includegraphics[width=\textwidth]{../images/02.png}
    \captionof{figure}{Настройка trunk на msk-ansusenko-sw-2.}
    \label{02}
\end{center}

\item На коммутаторе \texttt{msk-ansusenko-sw-3} настроили trunk на порту \texttt{g0/1} (Рис. \ref{03}):
  \begin{minted}{cisco}
    enable
    configure terminal
    interface g0/1
    switchport mode trunk
    end
    copy running-config startup-config
  \end{minted}

\begin{center}
    \centering
    \includegraphics[width=\textwidth]{../images/03.png}
    \captionof{figure}{Настройка trunk на msk-ansusenko-sw-3.}
    \label{03}
\end{center}

\item На коммутаторе доступа \texttt{msk-ansusenko-sw-1} настроили trunk на порту \texttt{FastEthernet0/24} (Рис. \ref{04}):
  \begin{minted}{cisco}
    enable
    configure terminal
    interface FastEthernet0/24
    switchport mode trunk
    end
    copy running-config startup-config
  \end{minted}

\begin{center}
    \centering
    \includegraphics[width=\textwidth]{../images/04.png}
    \captionof{figure}{Настройка trunk на msk-ansusenko-sw-1.}
    \label{04}
\end{center}

\end{enumerate}

\subsection{Настройка VTP-сервера и создание VLAN}

На коммутаторе \texttt{msk-ansusenko-sw-01}, выполняющем роль центрального, настроили VTP-сервер и создали необходимые VLAN (Рис. \ref{05}).

\begin{minted}{cisco}
enable
configure terminal
vtp mode server
vtp domain ansusenko
vtp password cisco
vlan 2
name management
vlan 3
name servers
vlan 101
name dk
vlan 102
name departaments
vlan 103
name adm
vlan 104
name other
end
\end{minted}

\begin{center}
    \centering
    \includegraphics[width=\textwidth]{../images/05.png}
    \captionof{figure}{Настройка VTP-сервера и создание VLAN на msk-ansusenko-sw-01.}
    \label{05}
\end{center}

\subsection{Настройка VTP-клиентов и портов доступа}

Коммутаторы \texttt{msk-ansusenko-sw-2}, \texttt{sw-3}, \texttt{sw-4} и \texttt{sw-1} настроили как VTP-клиенты. На интерфейсах указали принадлежность к соответствующим VLAN.

\begin{enumerate}
\item Настройка VTP-клиента на \texttt{msk-ansusenko-sw-2} (Рис. \ref{06}):
  \begin{minted}{cisco}
    configure terminal
    vtp mode client
    vtp domain ansusenko
    vtp password cisco
  \end{minted}

\begin{center}
    \centering
    \includegraphics[width=\textwidth]{../images/06.png}
    \captionof{figure}{Настройка VTP-клиента на msk-ansusenko-sw-2.}
    \label{06}
\end{center}

\item Настройка VTP-клиента на \texttt{msk-ansusenko-sw-3} (Рис. \ref{07}):
  \begin{minted}{cisco}
    configure terminal
    vtp mode client
    vtp domain ansusenko
    vtp password cisco
  \end{minted}

\begin{center}
    \centering
    \includegraphics[width=\textwidth]{../images/07.png}
    \captionof{figure}{Настройка VTP-клиента на msk-ansusenko-sw-3.}
    \label{07}
\end{center}

\item Настройка VTP-клиента и назначение портов доступа на \texttt{msk-ansusenko-sw-1} (Рис. \ref{08}):
  \begin{minted}{cisco}
    configure terminal
    vtp mode client
    vtp domain ansusenko
    vtp password cisco
    interface FastEthernet0/1
    switchport mode access
    switchport access vlan 101
    interface FastEthernet0/20
    switchport mode access
    switchport access vlan 104
  \end{minted}

\begin{center}
    \centering
    \includegraphics[width=\textwidth]{../images/08.png}
    \captionof{figure}{Настройка VTP-клиента и портов доступа на msk-ansusenko-sw-1.}
    \label{08}
\end{center}

\item Настройка VTP-клиента и диапазона портов на \texttt{msk-ansusenko-sw-4} в соответствии с примером из методички (п. 5.3.3) (Рис. \ref{09}):
  \begin{minted}{cisco}
    configure terminal
    vtp mode client
    interface range f0/1-5
    switchport mode access
    switchport access vlan 101
  \end{minted}

\begin{center}
    \centering
    \includegraphics[width=\textwidth]{../images/09.png}
    \captionof{figure}{Настройка VTP-клиента и портов доступа на msk-ansusenko-sw-4.}
    \label{09}
\end{center}

\end{enumerate}

\subsection{Проверка работы VTP и конфигурации VLAN}

\begin{enumerate}
\item Проверка статуса VTP на сервере \texttt{msk-ansusenko-sw-01} (Рис. \ref{10}):
  \begin{minted}{cisco}
    show vtp status
    show vlan brief
  \end{minted}

\begin{center}
    \centering
    \includegraphics[width=\textwidth]{../images/10.png}
    \captionof{figure}{Проверка VTP и VLAN на сервере.}
    \label{10}
\end{center}

Параметры VTP:
\begin{itemize}
  \item Operating Mode: Server
  \item Domain: ansusenko
  \item Configuration Revision: 13
  \item Number of existing VLANs: 11
\end{itemize}

\item Проверка статуса VTP на клиенте \texttt{msk-ansusenko-sw-1} (Рис. \ref{11}):
  \begin{minted}{cisco}
    show vtp status
    show vlan brief
  \end{minted}

\begin{center}
    \centering
    \includegraphics[width=\textwidth]{../images/11.png}
    \captionof{figure}{Проверка VTP и VLAN на msk-ansusenko-sw-1.}
    \label{11}
\end{center}

Параметры VTP:
\begin{itemize}
  \item Operating Mode: Client
  \item Configuration Revision: 13 (синхронизирован с сервером)
  \item Number of existing VLANs: 11
\end{itemize}
Назначения портов: Fa0/1 – VLAN 101, Fa0/20 – VLAN 104.

\item Проверка статуса VTP на клиенте \texttt{msk-ansusenko-sw-2} (Рис. \ref{12}):
  \begin{minted}{cisco}
    show vtp status
    show vlan brief
  \end{minted}

\begin{center}
    \centering
    \includegraphics[width=\textwidth]{../images/12.png}
    \captionof{figure}{Проверка VTP и VLAN на msk-ansusenko-sw-2.}
    \label{12}
\end{center}

\item Проверка статуса VTP на клиенте \texttt{msk-ansusenko-sw-3} (Рис. \ref{13}):
  \begin{minted}{cisco}
    show vtp status
    show vlan brief
  \end{minted}

\begin{center}
    \centering
    \includegraphics[width=\textwidth]{../images/13.png}
    \captionof{figure}{Проверка VTP и VLAN на msk-ansusenko-sw-3.}
    \label{13}
\end{center}

Назначения портов: Fa0/11 – VLAN 103.

\item Начальное состояние VTP на \texttt{msk-ansusenko-sw-4} (до синхронизации) (Рис. \ref{14}):
  \begin{minted}{cisco}
    show vtp status
    show vlan brief
  \end{minted}

\begin{center}
    \centering
    \includegraphics[width=\textwidth]{../images/14.png}
    \captionof{figure}{Начальное состояние VTP на msk-ansusenko-sw-4 (Revision = 0).}
    \label{14}
\end{center}

Параметры VTP:
\begin{itemize}
  \item Configuration Revision: 0
  \item Number of existing VLANs: 5 (только стандартные)
\end{itemize}

\end{enumerate}

\subsection{Настройка IP-адресации и проверка доступности}

\begin{enumerate}
\item Проверка связности управляющего VLAN 2 с сервера (Рис. \ref{15}):
  \begin{minted}{cisco}
    ping 10.128.1.3
    ping 10.128.1.4
    ping 10.128.1.6
  \end{minted}

\begin{center}
    \centering
    \includegraphics[width=\textwidth]{../images/15.png}
    \captionof{figure}{Проверка доступности коммутаторов в VLAN 2 с сервера.}
    \label{15}
\end{center}


\item Проверка связности с клиента \texttt{msk-ansusenko-sw-4} после синхронизации VTP (Рис. \ref{16}):
  \begin{minted}{cisco}
    ping 10.128.1.3
    ping 10.128.1.4
    ping 10.128.1.6
  \end{minted}

\begin{center}
    \centering
    \includegraphics[width=\textwidth]{../images/16.png}
    \captionof{figure}{Проверка доступности коммутаторов в VLAN 2 с msk-ansusenko-sw-4.}
    \label{16}
\end{center}

\item Проверка связности с клиента \texttt{msk-ansusenko-sw-1} (Рис. \ref{17}):
  \begin{minted}{cisco}
    ping 10.128.1.3
    ping 10.128.1.4
  \end{minted}

\begin{center}
    \centering
    \includegraphics[width=\textwidth]{../images/17.png}
    \captionof{figure}{Проверка доступности коммутаторов в VLAN 2 с msk-ansusenko-sw-1.}
    \label{17}
\end{center}

\item Проверка доступности устройств в одном VLAN и разных VLAN (Рис. \ref{18}):
  \begin{minted}{bash}
    ping 10.128.101.10
    ping 10.128.103.10
  \end{minted}

\begin{center}
    \centering
    \includegraphics[width=\textwidth]{../images/18.png}
    \captionof{figure}{Проверка ping между устройствами.}
    \label{18}
\end{center}
\item Выполним в симуляции движение к 4 коммутатору, проследим за тем, как и куда движутся пакеты (Рис. \ref{19}):

\begin{center}
    \centering
    \includegraphics[width=\textwidth]{../images/19.png}
    \captionof{figure}{Доставка пакетов}
    \label{19}
\end{center}

\end{enumerate}

\subsection{Изучение процесса передачи ICMP-пакета в режиме симуляции}

В режиме симуляции Packet Tracer изучили процесс передвижения ICMP-пакета по сети (Рис. \ref{20}).

\begin{enumerate}
\item В списке событий симуляции зафиксированы пакеты различных протоколов:
  \begin{itemize}
    \item \textbf{ICMP} - пинг между узлами (dk-ansusenko-01).
    \item \textbf{STP} - служебный трафик протокола Spanning Tree.
    \item \textbf{CDP} - протокол Cisco Discovery Protocol.
  \end{itemize}

\item При анализе ICMP-пакета (клик по конверту PDU) были изучены заголовки:
  \begin{itemize}
    \item \textbf{Канальный уровень (Ethernet II):} MAC-адреса отправителя и получателя, тип протокола 0x0800 (IPv4).
    \item \textbf{Сетевой уровень (IPv4):} IP-адреса отправителя и получателя, протокол 1 (ICMP), TTL.
    \item \textbf{Транспортный/прикладной уровень (ICMP):} тип сообщения (8 – echo request, 0 – echo reply), код, контрольная сумма, идентификатор, номер последовательности.
  \end{itemize}
\end{enumerate}

\begin{center}
    \centering
    \includegraphics[width=\textwidth]{../images/20.png}
    \captionof{figure}{События симуляции в Packet Tracer.}
    \label{20}
\end{center}

\section{Выводы}
В результате выполнения лабораторной работы были приобретены практические навыки по настройке VLAN на коммутаторах сети с использованием протоколов VTP и Trunk.

\end{document}
